\documentclass[11pt]{article}

\usepackage[utf8]{inputenc}
\usepackage[T1]{fontenc}
\usepackage{amsmath,amssymb}
\usepackage[a4paper,margin=25mm]{geometry}
\usepackage{xcolor}
\usepackage{hyperref}
\usepackage{enumitem}
\usepackage{microtype}
\usepackage{fancyhdr}
\microtypesetup{expansion=false}

\definecolor{assistantblue}{RGB}{32,82,147}
\definecolor{usergreen}{RGB}{31,122,86}
\definecolor{rulegray}{RGB}{205,210,216}

\hypersetup{
  colorlinks=true,
  linkcolor=assistantblue,
  urlcolor=assistantblue,
  pdfauthor={Conversation transcript},
  pdftitle={Closed-form rectangular CG overlaps}
}

\pagestyle{fancy}
\fancyhf{}
\fancyhead[L]{Closed-form rectangular CG overlaps}
\fancyhead[R]{Conversation transcript}
\fancyfoot[C]{\thepage}
\setlength{\headheight}{14pt}

\setlength{\parindent}{0pt}
\setlength{\parskip}{0.65em}
\setlist[itemize]{topsep=0.25em,itemsep=0.2em,leftmargin=1.6em}

\newcommand{\speaker}[2]{%
  \par\vspace{0.7em}%
  {\large\bfseries\color{#2}#1}\par
  \vspace{0.15em}\color{rulegray}\hrule\color{black}\vspace{0.55em}%
}

\newcommand{\Assistant}{\speaker{Assistant}{assistantblue}}
\newcommand{\User}{\speaker{User}{usergreen}}

\begin{document}

\begin{center}
  {\LARGE\bfseries Conversation transcript}\par
  \vspace{0.35em}
  {\Large Closed-form rectangular CG overlaps}\par
\end{center}

\vspace{0.8em}

\emph{This transcript begins at the requested sentence. It contains the visible user and
assistant messages only; internal tool calls and hidden instructions are omitted.}

\emph{Notation for external readers: throughout this transcript, the abbreviated bra is
defined by
\[
  \langle a,b,c,d\mid
  \equiv \langle a,b,c,d,0,0,0,0\mid
  := \langle \mathrm{HWS}\mid
  \bigl(e_{13}^{(1)}\bigr)^a
  \bigl(e_{14}^{(1)}\bigr)^b
  \bigl(e_{23}^{(1)}\bigr)^c
  \bigl(e_{24}^{(1)}\bigr)^d .
\]
Here the superscript \((1)\) means that every generator acts on the first tensor-product
site only; the four trailing zeros indicate that no generator acts on the second site.}

\Assistant

The exact small cases expose a much simpler scaling law than the matrix implementation:
every allowed output is
\[
  \mu=(s+p-a,\ s+p-b,\ b,\ a),
  \qquad 0\leq a\leq b\leq \min(s,p),
\]
and its complete overlap sector has only \(a+1\) entries. Their projective values are
\[
  \frac{\langle r,a-r,b-r,r\mid\Psi_\mu\rangle}
       {\langle0,a,b,0\mid\Psi_\mu\rangle}
  =\frac{(-1)^r}{\binom{b}{r}}.
\]
I've checked this exactly against the existing GT-matrix reconstruction through
\texttt{s=3,p=2}. The corresponding exponential overlap collapses to one Pl\"ucker power,
so extraction is linear in \texttt{min(s,p)}, not in representation dimension.

\User

you crashed my pc, whats the conclusion?

\Assistant

Sorry --- the live Wolfbook editing bridge accumulated stuck processes while I was updating
the notebook. I stopped them, but I should have backed off sooner. The workbook was not
successfully updated.

The mathematical conclusion is clean:

For
\[
  [s,s,0,0]\otimes[p,p,0,0],
\]
every output irrep is uniquely parameterized by
\[
  \mu=(s+p-a,\ s+p-b,\ b,\ a),
  \qquad 0\leq a\leq b\leq\min(s,p).
\]

Only these first-site bras can overlap:
\[
  (n_{13},n_{14},n_{23},n_{24})
  =(r,a-r,b-r,r),\qquad r=0,\ldots,a.
\]

After fixing
\[
  \langle0,a,b,0\mid\Psi_\mu\rangle=1,
\]
all overlaps are
\[
  \boxed{
    \langle r,a-r,b-r,r\mid\Psi_\mu\rangle
    =\frac{(-1)^r}{\binom{b}{r}}
  }.
\]

Equivalently, the complete exponential generating function is
\[
  \boxed{
    \frac{u_{23}^{\,b-a}}{a!\,b!}
    \left(u_{14}u_{23}-u_{13}u_{24}\right)^a
  }.
\]

All other
\(\langle n_{13},n_{14},n_{23},n_{24},0,0,0,0\mid\Psi_\mu\rangle\)
vanish by Cartan weight.

I verified this exactly against the existing explicit GT/tensor reconstruction through
\(s=3,p=2\), including:
\begin{itemize}
  \item \(a=b=1:\ \{1,-1\}\),
  \item \(a=1,b=2:\ \{1,-\tfrac12\}\),
  \item \(a=b=2:\ \{1,-\tfrac12,1\}\).
\end{itemize}

Thus the overlap extraction itself is representation-free and costs only \(O(a)\). The
Q-functions and bilinear-residue determinant can still be computed alongside it as Q-side
verification, without constructing \texttt{dim[rep]}, generator matrices, Kronecker products,
or \texttt{Psi}.

What remains is only safely placing this API into the workbook and running the requested
clean-kernel \(s=p=5\) timing. I did not complete that benchmark, so I cannot honestly claim
the measured sub-second result yet.

\User

Wait so youre saying there is an exact closed form expression for the CG coefficients?

\Assistant

Yes---but specifically for the restricted first-site overlaps you asked for, not automatically
for every CG coefficient in a conventional GT tensor basis.

For an output irrep
\[
  \mu=(s+p-a,\ s+p-b,\ b,\ a),
  \qquad 0\leq a\leq b\leq\min(s,p),
\]
the only Cartan-compatible overlaps are
\[
  \langle r,a-r,b-r,r,0,0,0,0\mid\Psi_\mu\rangle,
  \qquad r=0,\ldots,a.
\]

Up to one overall normalization, they have the exact closed form
\[
  \frac{\langle r,a-r,b-r,r\mid\Psi_\mu\rangle}
       {\langle0,a,b,0\mid\Psi_\mu\rangle}
  =\frac{(-1)^r}{\binom{b}{r}}.
\]

So their generating function is simply
\[
  \frac{u_{23}^{b-a}}{a!\,b!}
  (u_{14}u_{23}-u_{13}u_{24})^a
\]
when \(\langle0,a,b,0\mid\Psi_\mu\rangle=1\).

This completely determines the HWS projectively because the output irrep occurs with
multiplicity one. However:
\begin{itemize}
  \item it does not fix the absolute normalization;
  \item it is not yet a closed formula for every component of the HWS in Matteo's/GT tensor basis;
  \item the Q-determinant should supply/check the remaining normalization, but that identification
        still needs to be established carefully.
\end{itemize}

So the surprising result is: the entire family of overlaps you requested appears to have a
universal exact closed form, while the bilinear-residue machinery is mainly needed for the
Q-normalization and comparison with Matteo's convention.

\User

Is this formula known anywhere in the literature?

\Assistant

I'll check whether this precise binomial overlap formula is published explicitly, and distinguish
that from the well-known ingredients behind it---rectangular Littlewood--Richardson decomposition,
Pl\"ucker relations, and standard monomial theory.

\Assistant

The underlying formula is classical, but I have not found this exact four-index overlap formula
stated verbatim as an SU(4) Clebsch--Gordan formula.

It is a direct consequence of the Pl\"ucker/standard-monomial realization:
\begin{itemize}
  \item \([s,s,0,0]\) is the degree-\(s\) component of the homogeneous coordinate ring of
        \(\mathrm{Gr}(2,4)\).
  \item \(\mathrm{Gr}(2,4)\) has one quadratic Pl\"ucker relation.
  \item Products of rectangular Schur representations are multiplicity-free, so each permitted
        highest-weight vector is unique up to scale. This rectangular multiplicity-free result is
        classical; see Stembridge's
        \href{https://deepblue.lib.umich.edu/handle/2027.42/41839}
        {Multiplicity-Free Products of Schur Functions}.
  \item Standard monomial theory constructs these representations and their bases using Pl\"ucker
        coordinates; see Seshadri's
        \href{https://link.springer.com/book/10.1007/978-981-10-1813-8}
        {Introduction to the Theory of Standard Monomials}.
\end{itemize}

In our variables, the relevant standard monomial is simply
\[
  u_{23}^{\,b-a}
  \det\begin{pmatrix}
    u_{14}&u_{13}\\
    u_{24}&u_{23}
  \end{pmatrix}^{a}
  =u_{23}^{\,b-a}(u_{14}u_{23}-u_{13}u_{24})^a.
\]

Expanding it and accounting for the factorials in the exponential basis gives
\[
  \frac{\langle r,a-r,b-r,r\mid\Psi_\mu\rangle}
       {\langle0,a,b,0\mid\Psi_\mu\rangle}
  =(-1)^r\frac{r!(b-r)!}{b!}
  =\frac{(-1)^r}{\binom{b}{r}}.
\]

So I would describe it as:
\begin{quote}
  a simple specialization of the classical Pl\"ucker/standard-monomial highest-weight
  construction, expressed in our unusual raising-operator overlap convention.
\end{quote}

General SU(\(N\)) CG literature usually presents algorithms or GT-basis coefficients rather than
this restricted coherent-bra slice---for example the
\href{https://arxiv.org/abs/1009.0437}{Alex--Kalus--Huckleberry--von Delft algorithm}.
I did not locate a paper writing precisely
\[
  \langle r,a-r,b-r,r\mid\Psi\rangle
  =\frac{(-1)^r}{\binom{b}{r}}
\]
with our normalization.

Thus the closed form itself is probably not a fundamentally new representation-theory theorem.
What may be new and relevant to this project is showing that Matteo's bilinear Q-residue
determinant selects and normalizes exactly this standard-monomial state directly from the Gaudin
Q-functions.


\end{document}
